Electoral competitions can be viewed as situations in which each voter $i$ chooses among $P$ possible candidates, so that the response variable $Y_i$ is a $P$-dimensional vector with $Y_i \in \{0,1\}^P$.

However, in both the synthetic example and the application to U.S. electoral data, I will treat $Y_i$ as a binary choice. In the synthetic example, this is a consequence of the referendum rule, where only two options are available (vote Yes or vote No). In the empirical application, it follows from the American two-party system, where two parties dominate the political landscape and the remaining parties are not relevant for the analysis.

Hence, we can focus on one of the two possible votes and denote $Y_i$ by a binary indicator $Y_i \in \{0,1\}$, which takes value $1$ if the vote is cast in favor of the reference choice and $0$ otherwise.

Electoral precincts are the geographical units in which votes are aggregated, and they are the units of observation in the analysis. Therefore, we denote by $g$ a generic electoral precinct and by $G$ the number of electoral precincts.

As a consequence, the observed variable $\overline{Y}_g$ is a scalar with $\overline{Y}_g \in [0,1]$ that represents the share of the reference choice in precinct $g$.

We then consider a partition of individuals into $K$ mutually exclusive and exhaustive categories, which are the categories of the variable $X_i \in \{0,1\}^K$. Examples of such partitions could be the partition of voters into mutually exclusive and exhaustive income classes, or the partition of voters into mutually exclusive and exhaustive ethnic groups.

Finally, each electoral precinct $g$ is characterized by a vector of covariates $Z_g \in \mathbb{R}^d$, which may either describe geographic features (e.g., the location of the precinct in a particular region) or arise as aggregates of individual-level covariates (e.g., the education level of the voters in that precinct).

It is now necessary to introduce the assumptions that allow us to identify the model and estimate the parameters of interest.

The positivity assumption depends on the choice of covariates $Z_g$ and on the absence of collinearity between them and $\overline{X}_g$, so that the model is identifiable.

The bounded N assumption can be verified by verifying that no electoral precincts are too large.

Finally, we can impose a Coarsening At Random (CAR) assumption through an individual-level condition, which, following \citet{mccartan2025identificationsemiparametricestimationconditional}, we may call CAR-IND.

CAR-IND states that the individual-level response variable $Y_i$ is independent of the geography $g$, conditional on the individual-level covariates $X_i$ and the geography-level covariates $Z_g$. Formally, this can be expressed as follows:
\[
E[Y_i \mid G_i = g, X_i = x_i, Z_{G_i} = z_g] = E[Y_i \mid X_i = x_i, Z_{G_i} = z_g].
\]

From a practical point of view, this assumption implies that the covariates capture all the effects of geography on the individual-level response variable $Y_i$. In the electoral example, this means that White voters in precinct $g$ with characteristics $Z_g$ have the same probability of voting for the segregationist candidate as White voters in a different precinct $g'$ with the same characteristics $Z_{g'} = Z_g$.

Moreover, because $f(Z_g)$ is not necessarily injective, it is possible that two different electoral precincts $g$ and $g'$, even with different characteristics $Z_g \neq Z_{g'}$, have the same value $f(Z_g) = f(Z_{g'})$, and therefore the same expected value $\beta_g = \beta_{g'}$.

\citet{mccartan2025identificationsemiparametricestimationconditional} show that CAR-IND implies CAR.

Therefore, to perform semiparametric estimation of the model, we need to check bounded $N$ and then select covariates $Z_g$ that satisfy the testable positivity assumption and the untestable CAR-IND assumption.